% =====================================================================
\section{Atención por Caminos}\label{sec:model}
% =====================================================================

Definimos ahora la arquitectura. La observación rectora es que el
operador de caminos multi-salto de la Sección~\ref{sec:algebra} es estructuralmente una
\emph{atención}: agrega, para cada nodo objetivo, una combinación ponderada de
las características de los nodos que lo alcanzan. La autoatención del Transformer
\cite{vaswani2017attention} tiene la misma forma pero con soporte sobre todos los pares y
pesos \emph{aprendidos}; la atención sobre grafos \cite{velickovic2017graph} restringe el
soporte a un salto. La Atención por Caminos conserva los pesos aprendidos de la atención y
toma el soporte del álgebra de caminos: la alcanzabilidad por caminos multi-salto y dispersa
$\mathcal{S}_k$. En los términos de la Sección~\ref{sec:ground}, es la
actualización de nodo $\varphi$ aprendida del sustrato común---una representación
graduada de $Q$ cuyas aplicaciones de flecha se optimizan sobre el sistema
fundamental de vecindades.

% ---------------------------------------------------------------------
\subsection{La capa de Atención por Caminos}
% ---------------------------------------------------------------------

Fijemos un rango $k$ con soporte de alcanzabilidad $\mathcal{S}_k$
(Definición~\ref{def:walkop}). Una capa de Atención por Caminos con $H$ cabezas y dimensión
de cabeza $C$ transforma las características de los nodos $x\in\RR^{n\times d_{\text{in}}}$
así. Para cada cabeza $h$ calcula queries, keys y values mediante aplicaciones
lineales,
\[
  Q^h = x\,W^h_Q,\qquad K^h = x\,W^h_K,\qquad V^h = x\,W^h_V,
  \qquad W^h_\bullet\in\RR^{d_{\text{in}}\times C},
\]
y luego atiende \emph{sólo sobre los pares alcanzables}: para todo
$(j,i)\in\mathcal{S}_k$,
\begin{align}
  \ell^h_{ji} &= \frac{1}{\sqrt{C}}\,\langle Q^h_j,\,K^h_i\rangle, \label{eq:logit}\\
  \alpha^h_{ji} &= \softmax_{\,i:\,(j,i)\in\mathcal{S}_k}\bigl(\ell^h_{ji}\bigr),
    \label{eq:softmax}\\
  z^h_j &= \sum_{i:\,(j,i)\in\mathcal{S}_k}\alpha^h_{ji}\,V^h_i. \label{eq:agg}
\end{align}
El softmax en \eqref{eq:softmax} se toma, para cada objetivo $j$, exactamente sobre las
fuentes $i$ que alcanzan a $j$ por un camino de longitud $k$---un softmax
\emph{segmentado} sobre el soporte, nunca sobre los $n$ nodos. Las salidas de las cabezas se
concatenan (o promedian en la capa final) y se añade un término ``raíz'' residual:
\[
  \mathrm{WA}_k(x)_j \;=\; \bigl\Vert_{h=1}^{H} z^h_j \;+\; x_j W_{\!\text{root}}.
\]
Como \eqref{eq:logit}--\eqref{eq:agg} sólo tocan los pares de
$\mathcal{S}_k$, la capa cuesta $O(|\mathcal{S}_k|\,H\,C)$ en tiempo y memoria, no
$O(n^2)$. El soporte se precalcula una vez por grafo según la
Proposición~\ref{prop:count} (el patrón de no ceros de $\Adj^{\,k}$).
La Figura~\ref{fig:layer} resume el flujo de datos.

\begin{figure}[t]
\centering
\begin{tikzpicture}[
  >={Stealth[round]}, node distance=6mm,
  box/.style={draw,rounded corners=2pt,minimum height=7mm,inner sep=3pt,font=\small,fill=blue!4},
  alg/.style={draw,rounded corners=2pt,minimum height=7mm,inner sep=3pt,font=\small,fill=green!6},
  op/.style={font=\small},
  lbl/.style={font=\scriptsize,gray}]
  \node[box] (x) {$x$};
  \node[box,right=8mm of x] (q) {$Q=xW_Q$};
  \node[box,below=2mm of q] (k) {$K=xW_K$};
  \node[box,below=2mm of k] (v) {$V=xW_V$};
  \node[alg,right=10mm of k] (att)
    {$\displaystyle \alpha_{ji}=\softmax_{i:(j,i)\in\mathcal{S}_k}\!\frac{\langle Q_j,K_i\rangle}{\sqrt C}$};
  \node[box,right=10mm of att] (agg) {$z_j=\sum_i \alpha_{ji}V_i$};
  \node[box,right=8mm of agg] (out) {$\mathrm{WA}_k(x)$};
  % support callout
  \node[alg,below=8mm of att,fill=green!10] (supp)
    {soporte $\mathcal{S}_k=\{(j,i):(\Adj^{\,k})_{ij}>0\}$};
  \draw[->] (x) -- (q.west);
  \draw[->] (x) |- (k.west);
  \draw[->] (x) |- (v.west);
  \draw[->] (q.east) -- (att.north west);
  \draw[->] (k.east) -- (att.west);
  \draw[->] (att.east) -- (agg.west);
  \draw[->] (v.east) -- ([yshift=-2mm]agg.west);
  \draw[->] (agg) -- (out);
  \draw[->,green!50!black] (supp.north) -- (att.south)
    node[midway,right,lbl,text=green!50!black] {restringe los pares};
  \draw[->,rounded corners] (x.north) -- ++(0,6mm) -| (out.north)
    node[pos=0.25,above,lbl] {raíz (residual)};
\end{tikzpicture}
\caption{Una capa de Atención por Caminos en el rango $k$. Las queries, keys y values son
aplicaciones lineales de las características de los nodos; la atención se calcula \emph{sólo}
sobre los pares alcanzables por caminos $\mathcal{S}_k$ que proporciona el álgebra de caminos
(el patrón de no ceros de $\Adj^{\,k}$), y luego se agregan y se suman a un término raíz
residual. La caja verde de soporte es lo que distingue a la Atención por Caminos de la
atención densa ($\mathcal{S}_k = $ todos los pares) y de un salto ($\mathcal{S}_k = $ aristas).}
\label{fig:layer}
\end{figure}

\begin{remark}[Por qué importan los pesos aprendidos]
El operador de pesos fijos $W^{(k)}$ de la Definición~\ref{def:walkop} es el caso
particular en el que $\alpha^h_{ji}$ se reemplaza por la constante $n_k(i,j)/Z_k$.
Alcanza toda fuente de rango $k$ pero las pondera sólo por multiplicidad, sin poder
\emph{seleccionar} qué fuente porta la señal relevante. La atención aprendida
\eqref{eq:logit}--\eqref{eq:softmax} restaura esa selectividad
conservando el mismo soporte multi-salto. La Sección~\ref{sec:exp} muestra que la distinción
es decisiva.
\end{remark}

% ---------------------------------------------------------------------
\subsection{La red}
% ---------------------------------------------------------------------

Una red de Atención por Caminos de profundidad $L$ apila una capa por rango
$k=1,\dots,L$. La capa $k$ usa el soporte $\mathcal{S}_k$ derivado de
$\Adj^{\,k}$, así que las capas sucesivas leen caminos progresivamente más largos. Cada
capa va seguida de una normalización de capa, una no linealidad \textsc{elu} y
dropout; un perceptrón multicapa final lleva la incrustación de nodos a las salidas de la tarea:
\[
  x^{(k)} \;=\; \mathrm{dropout}\Bigl(\mathrm{elu}\bigl(\mathrm{LN}\bigl(
  \mathrm{WA}_k\bigl(x^{(k-1)}\bigr)\bigr)\bigr)\Bigr),\qquad
  \hat y \;=\; \mathrm{MLP}\bigl(x^{(L)}\bigr).
\]
La normalización de capa importa en la práctica: sin ella la atención aprendida tiene
alta varianza entre inicializaciones (alcanza el óptimo en algunas semillas y
en otras no); con ella, más un breve warmup y recorte de
gradiente, el entrenamiento converge de forma fiable (Sección~\ref{sec:exp}).
El Algoritmo~\ref{alg:wa} resume una capa.

\begin{algorithm}[t]
\caption{Capa de Atención por Caminos en el rango $k$ (dispersa).}\label{alg:wa}
\begin{algorithmic}[1]
\Require características $x\in\RR^{n\times d_{\text{in}}}$; soporte
  $\mathcal{S}_k=\{(j,i):(\Adj^{\,k})_{ij}>0\}$
\State $Q,K,V \gets xW_Q,\,xW_K,\,xW_V$ \Comment{por cabeza; formas $n\times HC$}
\For{each $(j,i)\in\mathcal{S}_k$ and head $h$}
  \State $\ell^h_{ji}\gets \langle Q^h_j,K^h_i\rangle/\sqrt{C}$
\EndFor
\State $\alpha^h_{ji}\gets$ $\softmax$ segmentado sobre $\{i:(j,i)\in\mathcal{S}_k\}$
\State $z^h_j\gets \sum_i \alpha^h_{ji}V^h_i$
  \Comment{scatter-add hacia los objetivos}
\State \Return $\bigl\Vert_h z^h \;+\; xW_{\!\text{root}}$
\end{algorithmic}
\end{algorithm}

% ---------------------------------------------------------------------
\subsection{Relación con el paso de mensajes y los Transformers}
% ---------------------------------------------------------------------

La Tabla~\ref{tab:relation} sitúa la Atención por Caminos entre los operadores relacionados según dos
ejes: el \emph{soporte} de interacción y si los pesos son
\emph{aprendidos}. Una capa de atención sobre grafos de un salto \cite{velickovic2017graph} aprende
pesos pero sólo sobre los vecinos directos, y debe propagar la señal de largo alcance
a través de $L$ capas apiladas---precisamente el régimen donde el over-squashing es más severo. Un
Transformer de grafo completo \cite{vaswani2017attention} aprende pesos sobre todos los pares,
pagando $O(n^2)$ y descartando la estructura del grafo salvo que se reinyecte mediante
codificaciones posicionales. El operador de caminos de pesos fijos tiene el soporte multi-salto
adecuado pero no puede seleccionar. La Atención por Caminos ocupa la celda restante: soporte
multi-salto y consciente de la estructura \emph{y} pesos aprendidos y selectivos.

\begin{table}[t]
\centering
\caption{La Atención por Caminos frente a operadores de agregación relacionados.}
\label{tab:relation}
\begin{tabular}{lccc}
\toprule
\textbf{Operador} & \textbf{Soporte} & \textbf{Pesos} & \textbf{Costo} \\
\midrule
GCN / paso de mensajes       & 1 salto      & fijos (grado) & $O(|Q_1|)$ \\
Atención sobre grafos (GAT)       & 1 salto      & aprendidos        & $O(|Q_1|)$ \\
Operador de caminos $W^{(k)}$     & alcanzable por caminos & fijos (multiplicidad) & $O(|\mathcal{S}_k|)$ \\
Transformer de grafo           & todos los pares  & aprendidos        & $O(n^2)$ \\
\textbf{Atención por Caminos (nuestra)} & \textbf{alcanzable por caminos} & \textbf{aprendidos} & $O(|\mathcal{S}_k|)$ \\
\bottomrule
\end{tabular}
\end{table}

% ---------------------------------------------------------------------
\subsection{Un tratamiento formal sobre la familia de cuellos de botella}
\label{sec:theory}
% ---------------------------------------------------------------------

Sobre la familia de cadenas con cuello de botella de la Sección~\ref{sec:exp},
la separación entre el paso de mensajes de un salto y Walk Attention se puede
\emph{demostrar}, no solo medir. Recordemos la construcción: $K$ vértices
fuente llevan una \emph{clave} one-hot---la asignación de claves a fuentes es
una permutación aleatoria uniforme $\pi$---y un \emph{valor} one-hot; siguen
$d$ capas de cuello de botella $L_1,\dots,L_d$ de $M$ vértices intercambiables
y sin atributos, con capas consecutivas totalmente conectadas; un vértice
objetivo lleva una \emph{consulta} one-hot $q$ y debe producir el valor de la
fuente cuya clave es $q$ (etiqueta $\pi^{-1}(q)$, azar $1/K$). En todo lo que
sigue las redes tienen la profundidad estándar $L=d+1$---una capa por salto,
como en el protocolo de \textsc{NeighborsMatch} \cite{alon2021bottleneck} y en
nuestros experimentos---y ancho oculto $w$.

\begin{proposition}[Colapso de capas]\label{prop:collapse}
Para toda red de paso de mensajes sobre este grafo cuya actualización
$h_v' = U\bigl(h_v,\mathrm{AGG}\{h_u : u\to v\}\bigr)$ usa los mismos $U$ y
$\mathrm{AGG}$ en cada vértice: tras cada ronda, los $M$ vértices de cada capa
de cuello de botella llevan estados idénticos. En consecuencia, toda la
influencia de las fuentes sobre el objetivo pasa por un único estado
$w$-dimensional por capa---mientras que el número de recorridos de longitud
$(d{+}1)$ que la transportan es $K M^{d}$ (Proposición~\ref{prop:count}).
\end{proposition}

\begin{proof}
Inducción sobre la ronda. Los vértices de una capa de cuello de botella
comienzan con atributos idénticos (cero), y dos vértices de la misma capa
tienen el mismo multiconjunto de vecinos de entrada (todas las fuentes para
$L_1$; todo $L_{r-1}$ para $L_r$), cuyos estados son iguales por inducción;
$U$ y $\mathrm{AGG}$ compartidos producen actualizaciones iguales. \qed
\end{proof}

\begin{proposition}[La primera agregación lineal es ciega a la permutación]
\label{prop:blind}
Supóngase además que los mensajes son lineales en el estado del emisor con
coeficientes independientes del contenido---como en GCN (mensajes lineales
normalizados por grado) y GIN (suma de los estados vecinos). Entonces, con
$L=d+1$ capas, la salida del objetivo depende de las fuentes solo a través de
$\sum_{s}x_s$, que es \emph{la misma para toda} $\pi$: las casillas de clave de
una permutación suman el vector de unos. La precisión esperada es por tanto
exactamente $1/K$, para toda profundidad $d$, todo ancho $w$ y toda elección de
parámetros.
\end{proposition}

\begin{proof}
Los vecinos de entrada del objetivo son $L_d$ (y él mismo); desenrollando, el
estado del objetivo tras la ronda $d{+}1$ depende de las fuentes solo a través
del estado de $L_1$ tras la ronda~$1$: la información de las fuentes que entra
a $L_1$ en la ronda $t$ necesita $d$ saltos más para llegar al objetivo, y
llega en la ronda $t+d\le d+1$ solo si $t\le1$. En la ronda $1$ el mensaje de
la fuente $s$ es una aplicación lineal fija del atributo crudo $x_s$, con un
único coeficiente compartido (las fuentes son isomorfas en el grafo), así que
el estado de $L_1$ es función de $\sum_s x_s$. Con claves $e_{\pi(s)}$, el
bloque de claves de $\sum_s x_s$ es $\sum_s e_{\pi(s)}=(1,\dots,1)$ para toda
$\pi$, y el bloque de valores es igualmente independiente de $\pi$; la salida
es por tanto independiente de $\pi$. Dada la consulta $q$, la etiqueta
$\pi^{-1}(q)$ es uniforme en $\{1,\dots,K\}$, de modo que toda predicción
independiente de $\pi$ tiene precisión esperada $1/K$. \qed
\end{proof}

\begin{proposition}[Una sola capa de Walk Attention basta]\label{prop:wa}
En el alcance $k=d{+}1$ el soporte del objetivo es exactamente el conjunto de
las $K$ fuentes (el grafo es graduado; Proposición~\ref{prop:fns}). Elíjase
$W_Q$ proyectando la casilla de consulta escalada por cualquier $\beta>0$,
$W_K$ la casilla de clave, $W_V$ la casilla de valor, y $W_{\!\text{root}}=0$.
Entonces la salida de Walk Attention en el objetivo es
$z_t=\sum_s \alpha_s\,e_{v(s)}$, cuya mayor entrada está en la coordenada del
valor correcto, en \emph{toda} instancia y para todos $d$ y $M$.
\end{proposition}

\begin{proof}
El logit de la fuente $s$ es $\beta\langle e_q, e_{\pi(s)}\rangle
=\beta\,\mathbf{1}[\pi(s)=q]$, así que el softmax sobre las $K$ fuentes
alcanzables pone peso $e^{\beta}/(e^{\beta}+K-1)$ en la única fuente
coincidente $s^\ast$ y pesos iguales y menores en el resto. Los valores
$e_{v(s)}$ son vectores de la base, distintos, de modo que la mayor coordenada
de $z_t$ es $\alpha_{s^\ast}$, en la coordenada $v(s^\ast)=\pi^{-1}(q)$, la
etiqueta. \qed
\end{proof}

\begin{remark}\label{rem:mechanisms}
Las proposiciones delimitan los mecanismos observados en la
Sección~\ref{sec:exp}. GAT escapa a la Proposición~\ref{prop:blind}---sus
coeficientes de la primera ronda dependen del contenido de la fuente---pero
aún debe hacer pasar la asociación clave--valor por el canal de un solo vector
de la Proposición~\ref{prop:collapse}, re-codificado $d{+}1$ veces:
empíricamente queda por encima del azar y se degrada con la profundidad. El
operador de caminos de pesos fijos lee las fuentes directamente pero con pesos
\emph{iguales} (todas las multiplicidades valen $M^{d}$), así que debe
incrustar el multiconjunto clave--valor completo en $w$ dimensiones en lugar
de seleccionar una fuente: se estanca entre ambos. Walk Attention lee
directamente \emph{y} selecciona (Proposición~\ref{prop:wa}), y resuelve la
tarea.
\end{remark}
